\section{Harnack inequalities and shrinking holomorphic charts}
\label{sec:charts}

The joint positivity of the cotangent norm will now give an exhausting
system of inverse charts with quantitative transition estimates.
We first prove positive-time completeness and construct uniform local
charts.  These charts may have several sheets, so continuing their
inverse germs across compact sets also requires a topological step:
we will contract loops before constructing single-valued inverse branches.
For the quantitative estimates, volume loss must force contraction in
every direction, with an exponent common to all compact sets.
The local chart and path-lifting constructions have antecedents in Chau--Tam
\cite[Section~2]{ChauTamSimply}\cite[Sections~2--3]{ChauTamStein}.
Throughout this Section, a metric on a cotangent space is the Hermitian
dual metric, including the transpose in its matrix representation.

\subsection{From joint positivity to completeness}

The differential Harnack inequalities for K\"ahler--Ricci flow have
their classical antecedent in Cao's work \cite{CaoHarnack}.
Here we derive the needed inequalities from joint positivity before
using them to prove completeness of the evolving metric.

\begin{proposition}\label{prop:harnack-completeness}
The flow of Theorem~\ref{thm:localized-flow} has nonnegative holomorphic
bisectional curvature and is complete at every positive time.  Its
Ricci tensor is positive definite.  Writing $A=\Ric(h_t)$ and
$R=\tr_{h_t}A$, one has, in fixed holomorphic coordinates,
\begin{align}
 A_t+A h_t^{-1}A+t^{-1}A&\ge0,
 \label{chart:matrix-harnack}\\
 R_t+t^{-1}R+2\operatorname{Re}\partial_XR+A(X,\overline X)&\ge0
 \qquad(X\in T^{1,0}M).
 \label{chart:scalar-harnack}
\end{align}
In particular $h_{t'}\le h_t$ for $t'\ge t$, $R(o,t)\le C$ for
$t\ge0$, and
\begin{equation}\label{chart:integrated-harnack}
 R(x,t)\le R(y,t')\frac{t'}t
 \exp\!\left(\frac{C\dist_{h_t}(x,y)^2}{t'-t}\right)
 \qquad(0<t<t').
\end{equation}
On each fixed compact set, $h_t\le C_Kt^{-a_K}g$ for $t\ge1$, with
$a_K>0$.
\end{proposition}

\begin{proof}
By Theorem~\ref{thm:joint-positivity}, the dual squared norm of the
tangent metric pulled back to $(x,w)$, where $t=e^{w+\bar w}$, is
plurisubharmonic.  Completing the square in the free first derivative
of a cotangent section identifies this assertion with Griffiths
nonnegativity of the pulled-back tangent metric.  At a spatially
K\"ahler-normal point its curvature blocks are
\[
 \Theta_{i\bar j\,\alpha\bar\beta}
   =\Rm_{i\bar j\alpha\bar\beta},\qquad
 \Theta_{w\bar j\,\alpha\bar\beta}
   =t\nabla_{\bar j}A_{\alpha\bar\beta},\qquad
 \Theta_{w\bar w}=tA+t^2(A_t+A h_t^{-1}A).
\]
The spatial block gives nonnegative bisectional curvature and hence
$A\ge0$ and metric monotonicity.  The last block gives
\eqref{chart:matrix-harnack}.  Trace the curvature in its fiber
indices and evaluate on the spacetime direction $(X,t^{-1})$.
The identities
\[
 \tr_{h_t}A_t=R_t-|A|_{h_t}^2,
 \qquad \tr_{h_t}(A h_t^{-1}A)=|A|_{h_t}^2,
 \qquad \tr_{h_t}\nabla_X A=\partial_XR
\]
give \eqref{chart:scalar-harnack}.

At a fixed point choose a moving unitary frame satisfying
$E_t=\tfrac12 A^{\#}E$, where $A^{\#}=h_t^{-1}A$.
Differentiation shows that
\[
 \frac{d}{dt}\bigl[tA(E,\overline E)\bigr]
   =t\bigl(A_t+A h_t^{-1}A+t^{-1}A\bigr)(E,\overline E)\ge0.
\]
Initial strict bisectional positivity and local smooth convergence
give $A>0$ at sufficiently small positive times on any prescribed
compact set.  The preceding inequality propagates this positivity
and gives $tA\ge a_Kh_t$ there for $t\ge1$, after changing the compact
constant.  Integrating $\partial_t h_t=-A$ proves the asserted
power contraction.  Applied to the trace, the same inequality gives
monotonicity of $tR(o,t)$.  The linear bound for $\rho(o,t)$ in
Theorem~\ref{thm:localized-flow} implies, for $t\ge1$,
\[
 tR(o,t)\log2
 \le\int_t^{2t}R(o,v)\,dv
 =\rho(o,2t)-\rho(o,t)\le C(1+2t).
\]
Smoothness on compact time intervals supplies the remaining
basepoint bound.

For clarity, the next integration precedes the completeness proof.
Choose a piecewise smooth path from $x$ to $y$ whose $h_t$ length is
within $\eps$ of their distance, and run it at constant $h_t$ speed
on $[t,t']$.  Such paths exist by the definition of the distance.
Along this path, \eqref{chart:scalar-harnack}, divided by $R>0$,
and $A\le Rh_v$ give
\[
 \frac{d}{dv}\log R(\gamma(v),v)
 \ge-\frac1v-C|\dot\gamma(v)|_{h_v}^2.
\]
Since $h_v\le h_t$, its energy is at most
$(\dist_{h_t}(x,y)+\eps)^2/(t'-t)$.  Integrating and then letting
$\eps\downarrow0$ proves \eqref{chart:integrated-harnack} without
using a minimizing geodesic or completeness of $h_t$.

Fix $T<\infty$.  Put $d=\dist_{h_t}(x,o)$ and apply
\eqref{chart:integrated-harnack} with $y=o$ and
$t'=t+1+d^2$.  The basepoint bound gives
\begin{equation}\label{chart:quadratic-curvature}
 R(x,t)\le\frac{C_T(1+d^2)}t,
 \qquad 0<t\le T.
\end{equation}
Let $D_L=B_g(o,L)$, and let $l_L(t)$ be the infimum of lengths of
paths from $o$ to first exit from $D_L$, measured in $h_t$.
The initial metric is complete, so $\overline D_L$ is compact.
There is a minimizing segment realizing $l_L(t)$, contained in
$\overline D_L$, and every initial part of it minimizes to its
endpoint.  On this segment the distance in
\eqref{chart:quadratic-curvature} is at most $l_L(t)$.

Write $l=l_L(t)$.  In the summed real index forms use perpendicular
parallel fields times a cutoff which rises from zero to one on
an endpoint interval of width
\[
 a=\min\left(\frac l2,\frac{\sqrt t}{1+l}\right).
\]
If $l\ge2\sqrt t/(1+l)$, the index inequality and the two endpoint
curvature bounds imply
\[
 \int_0^l\Ric_{h_t}(\dot\gamma,\dot\gamma)\,dr
 \le \frac C a+\frac{C_T(1+l^2)a}{t}
 \le\frac{C_T(1+l)}{\sqrt t}.
\]
Here and below fixed real-versus-complex normalization factors are
included in $C_T$.  If $l<2\sqrt t/(1+l)$, use
\eqref{chart:quadratic-curvature} on the whole segment to obtain the
same bound.  The function $l_L$ is locally Lipschitz in time, since
the metrics are smoothly equivalent on the compact closure.  At a
differentiability time, compare a current minimizing segment with
its length at a slightly earlier time.  The metric variation formula
therefore gives
\[
 l_L'(t)\ge-\frac{C_T(1+l_L(t))}{\sqrt t}
 \quad\text{for almost every }t\in(0,T].
\]
Consequently
\begin{equation}\label{chart:exit-distance}
 1+l_L(t)\ge(1+L)e^{-2C_T\sqrt t}.
\end{equation}
As $L\to\infty$, every divergent path has infinite $h_t$ length.
Equivalently, each finite $h_t$ ball is trapped in a compact initial
ball; its closure is compact by local equivalence of the smooth
metrics.  This proves completeness.
\end{proof}

\begin{lemma}[Uniform immersed charts]\label{chart:local-charts}
For every fixed compact $K\subset M$, there are $r_K,c_K,C_K>0$
such that, for sufficiently large $t$ and every $p\in K$, there is
a holomorphic local biholomorphism
\[
 \Phi:B_{\C^n}(r_K)\longrightarrow M,\qquad
 \Phi(0)=p,\qquad d\Phi(0)\text{ unitary for }h_t,
\]
with $c_Kg_{\rm eucl}\le\Phi^*h_t\le C_Kg_{\rm eucl}$.
The pullback metrics have bounds of every fixed derivative order
on smaller balls.  The corresponding local smooth compactness
statement holds for rescaled flows whenever curvature is uniformly
bounded on fixed-radius moving balls during a fixed two-sided time
interval.
\end{lemma}

\begin{proof}
For $t\ge1$, \eqref{chart:integrated-harnack}, with a fixed future
time increment, bounds $R$ on every fixed-radius $h_t$ ball about
$o$.  Nonnegative bisectional curvature gives $|\Rm|\le C_nR$.
Centers in $K$ are handled by a compact family of paths from $o$:
their $h_t$ lengths are bounded, so balls about these centers lie
in a bounded-radius ball about $o$.

To apply local curvature derivative estimates, take a ball for the
earliest time of a short time interval.  It remains inside the
corresponding larger later-time balls because metrics decrease.
The curvature bound gives metric equivalence on this fixed region
throughout the interval.  A path measured in the later metric
cannot first exit the earlier ball at arbitrarily small length;
thus a smaller later-time ball is contained in this region.
Local Ricci-flow derivative estimates
\cite{Shi}\cite[Theorem~1.4]{Kotschwar} now give all curvature
derivatives on smaller balls and inner time intervals.
Here the estimate is applied on shifted sufficiently short cylinders,
with a time buffer before each observation slice.  Metric equivalence
keeps the initial balls required by the local estimate compactly
inside the fixed region.  This construction applies verbatim after
a fixed rescaling.

Here is why no injectivity-radius hypothesis enters the chart
statement; compare \cite[Lemma~2.2]{ChauTamSimply}.
Pull back by the exponential map on a small tangent ball.  The
curvature bound excludes conjugate points at this scale, although
the exponential map may identify distinct points.  Jacobi-field
estimates and their differentiated equations give uniformly
controlled pullback metrics and parallel complex structures, with
a fixed Euclidean lower metric bound.  In particular any sequence
has a smoothly convergent subsequence on a smaller tangent ball.
The identities $J^2=-\Id$ and $N_J=0$ pass to the smooth limit.
The Newlander--Nirenberg theorem supplies local holomorphic coordinates
for the limit complex structure \cite{NewlanderNirenberg}
\cite[Theorem~1.1]{HillTaylor}.  In the smooth case their components
satisfy the smooth elliptic equation $\bar\partial_J^*\bar\partial_J f=0$,
so interior elliptic regularity makes these coordinates smooth.
Restrict them to a still smaller ball.  Their squared radius is uniformly strictly
plurisubharmonic for all sufficiently close complex structures.
Solve the small $\dbar$ errors of these coordinate functions on
a smaller strictly pseudoconvex coordinate ball.  One can use the
weighted $L^2$ estimate \cite{Hormander,Demailly}, viewing scalar
functions as top forms with values in the anticanonical bundle.
Give this bundle the metric induced by the original uniformly
controlled pullback metric.  A fixed large multiple of the squared
radius pays its uniformly
bounded curvature, and the strictly plurisubharmonic exhaustion
of the ball supplies an auxiliary complete K\"ahler metric.
The source errors tend to zero smoothly, so the correcting
functions tend to zero in local $L^2$.  We first justify smoothness of
each correction, then obtain uniform estimates on a fixed smaller
real ball.  Write the corrections as $u_j$
and their equations as $\dbar_{J_j}u_j=\eta_j$.
For each index separately, the existing complex structure $J_j$
has smooth holomorphic coordinates, with no uniform radius asserted.
In those coordinates,
\[
 \Delta_{\mathrm{eucl}}u_j
   =4\sum_a\partial_{z_a}\eta_{j,a}
\]
is an identity of distributions with smooth right side.  A compactly
supported smooth extension of the right side, convolved with the
Euclidean fundamental solution, gives a smooth particular solution
on a smaller ball.
The remaining harmonic distribution is smooth: its mollifications
satisfy the mean-value identity against a fixed smooth radial
unit-mass kernel, and passage to the distributional limit identifies
the distribution with its smooth convolution by that kernel.
Thus each $u_j$ is qualitatively smooth.

For uniform estimates return to the fixed real ball and use
$P_j=\dbar_{J_j}^*\dbar_{J_j}$ with the original controlled
pullback K\"ahler metric, rather than the auxiliary weighted metric.
The scalar K\"ahler identity identifies $P_j$, up to a fixed sign
and normalization, with the real Laplace--Beltrami operator.
Its coefficients have uniform bounds of every fixed order and a
uniform ellipticity lower bound.  Moreover
$P_ju_j=\dbar_{J_j}^*\eta_j=:F_j$ tends smoothly to zero.
Testing this equation with $\chi^2\overline{u_j}$ and absorbing
the cross term gives the Caccioppoli estimate
\[
 \int\chi^2|\nabla u_j|^2
 \le C\int_{\supp\chi}(|u_j|^2+|F_j|^2)
\]
for a fixed interior cutoff.  The qualitative smoothness just
proved legitimizes this test.  Uniform metric equivalence gives a
$W^{1,2}$ bound in the fixed real coordinates.  Interior elliptic
estimates \cite[Theorems~8.8 and~8.10]{GilbargTrudinger}, on
successively smaller fixed balls, now yield
\[
 \|u_j\|_{H^{m+2}(B')}
 \le C_m\bigl(\|u_j\|_{L^2(B)}+\|F_j\|_{H^m(B)}\bigr).
\]
Choose $m+2>k+n$ and use Sobolev embedding in real dimension $2n$
\cite[Section~5.6.3, Theorem~6]{EvansPDE}.  This proves convergence
to zero in every fixed $C^k$ norm on smaller balls.  The corrected
coordinate maps are therefore $C^1$ close to the limiting coordinate
map.  On a sufficiently small fixed ball their derivatives are
uniformly close to its invertible center derivative.  They are
injective on that ball, and their images contain a common ball about
their center values, by the quantitative inverse function theorem.
Subtract the coordinate center values and invert on this common ball.
Composing with the exponential parametrization and making a bounded
linear normalization, with one further fixed reduction of radius,
gives the required maps into $M$ with unitary center derivative.
If no common radius and bounds existed for the original family, a
sequence witnessing their failure would have a smoothly convergent
subsequence of pullbacks.  The construction just given supplies common
charts on that subsequence, a contradiction.  This proves the uniform
chart assertion.

For a flow, use the exponential parametrization at its central
time as a fixed parametrization.  The curvature bound controls
the time derivative of the metric; the first Ricci derivative
controls that of its connection.  Induction using the higher
covariant derivative estimates gives ordinary space-time
derivative bounds for the metrics and their parallel complex
structures.  Arzel\`a--Ascoli on smaller balls gives the claimed
local smooth limits.
\end{proof}

\subsection{Total distortion in the initial metric}

In an $h_t$-unitary chart at $x$, the product of the initial-metric
singular lengths is $e^{\rho(x,t)/2}$.  This determinant identity alone
allows the growth to occur in only some directions.  The next lemma
bounds the greatest singular length by a fixed power of the least,
so volume loss will force contraction in every direction.

\begin{lemma}[Static distortion]\label{lem:static-distortion}
Fix $r,c>0$ and a compact set $K\subset M$.  Suppose
$\Phi:B_{\C^n}(r)\to M$ is a holomorphic local biholomorphism,
$\Phi(0)\in K$, and $\widetilde g=\Phi^*g\ge c g_{\rm eucl}$.
Let $A$ and $B$ be the least and greatest singular lengths of
$d\Phi(0)$, using the Euclidean metric in the domain and $g$ in
the target.  There are constants depending only on $K,r,c$ and
the initial geometry on a compact neighborhood of $K$ such that
\begin{equation}\label{chart:static-distortion}
 B\le C(1+A)^C.
\end{equation}
In particular these constants are independent of the particular
chart and its total distortion.
\end{lemma}

\begin{proof}
We construct a plurisubharmonic weight whose upper bound is of order
$\log(1+A)$ and whose Levi form is strict near the center.
Its moving-ball definition will also give a nonnegative lower bound
on the support of the weighted $\dbar$ construction.  These bounds
make the extension estimate independent of the total distortion.
We then extend a local function whose differential detects $B$;
its $L^2$ bound gives the required polynomial bound.

Choose a constant holomorphic covector $\alpha=df$ with
$|\alpha|_{\rm eucl}=A$ and $|\alpha|_{\widetilde g}(0)=1$.
The logarithmic dual norm
\[
 \psi=\log|\alpha|_{\widetilde g}^2
 \quad\text{satisfies}\quad
 \psi(0)=0,\qquad \psi\le C+2\log A.
\]
For a tangent $X$, the curvature formula for this logarithmic norm is
\[
 \dd\psi(X,\overline X)
 =\frac{\Rm(X,\overline X,v,\overline v)}{|v|^2}
  +\frac{|D'_X\alpha|^2|\alpha|^2
       -|\langle D'_X\alpha,\alpha\rangle|^2}{|\alpha|^4},
\]
where $v$ is the metric dual of $\alpha$ and all norms in this
formula use $\widetilde g$.  The second summand is nonnegative.
Initial bisectional positivity therefore makes $\psi$
plurisubharmonic.  On any region whose target image
lies in a prescribed compact neighborhood $K^+$ of $K$, it gives
\begin{equation}\label{chart:strict-covector}
 \dd\psi\ge\kappa\widetilde g
\end{equation}
with $\kappa>0$ uniform.  The normalization of $\alpha$ cancels
from this estimate.

Shrink the Euclidean domain by a fixed factor.  Choose $r_0>0$
so small that any path of $\widetilde g$ length at most $3r_0$
starting in this smaller domain stays in $B(r)$; the lower
metric bound makes this choice uniform.  The small closed metric
balls used below are compact there and their distances are
attained by minimizing segments.  Put
\begin{equation}\label{chart:metric-ball-envelope}
 \Psi(z)=\max_{\dist_{\widetilde g}(z,y)\le r_0}\psi(y).
\end{equation}
This function is continuous.  Compactness gives upper
semicontinuity.  For lower semicontinuity, move a maximizing
endpoint a little towards the center along a minimizing segment;
the resulting endpoint is admissible for all nearby centers and
its value tends to the maximum.

We verify the Levi inequality for this moving-ball maximum.
Let $\varphi$ be a smooth upper test at $z$, let $y$ realize the
maximum, and prescribe a nonzero complex tangent $V$ at $z$.
Choose a holomorphic center germ with this derivative.  If
$\dist_{\widetilde g}(z,y)<r_0$, any holomorphic endpoint germ
with derivative equal to the parallel translate of $V$ along a
minimizing segment stays admissible for a sufficiently small
parameter.  Suppose instead that the constraint is active.
For two real parameter directions corresponding to $V$ and $JV$,
choose trial paths whose first variation fields are the parallel
translates along this segment.  Their first length variations
vanish.  The trace of their second length variations is
\[
 -\int_0^{r_0}
 \bigl[\Rm(T,V^{\parallel},V^{\parallel},T)
       +\Rm(T,JV^{\parallel},JV^{\parallel},T)\bigr]\,dr<0.
\]
Tangential components make no curvature contribution, and
holomorphic endpoint germs have zero trace covariant
acceleration, so there are no omitted trace boundary terms.
The strict sign is precisely the initial bisectional sign.

The trace inequality alone is insufficient to preserve a
two-parameter length constraint.  At the far endpoint we may
prescribe its holomorphic second jet freely while retaining
its first derivative.  Changing that jet by a real tangent $W$
(regarded as a complex tangent by $J$) changes the real length
Hessian by
\[
 \begin{pmatrix}
 \langle W,T\rangle&\langle JW,T\rangle\\
 \langle JW,T\rangle&-\langle W,T\rangle
 \end{pmatrix}
\]
at the final endpoint.  Vectors in the real span of $T,JT$
realize every trace-free symmetric matrix of size two.
Choose the jet to remove the trace-free part of the length
quadratic.  The trace remains
strictly negative; hence the full quadratic is negative definite.
Smooth trial paths joining these holomorphic endpoint germs
then have length less than $r_0$ for every sufficiently small
nonzero parameter.  The actual distance is no greater than this
trial length.  This proves admissibility also at an active
constraint, without differentiating a possibly nonsmooth
distance function.

On either kind of admissible endpoint germ $y(a)$,
$\varphi(z(a))\ge\Psi(z(a))\ge\psi(y(a))$, with equality at
$a=0$.  Taking the parameter Levi derivative proves
$\dd\Psi\ge0$ in upper-test sense, and thus in the
plurisubharmonic sense.  If
$\dist_{\widetilde g}(0,z)<r_0/4$, all the relevant paths and
endpoints map into a fixed compact $K^+$; applying
\eqref{chart:strict-covector} to the transported tangent gives
\begin{equation}\label{chart:strict-envelope}
 \dd\Psi\ge\kappa\widetilde g
 \quad\text{on }B_{\widetilde g}(0,r_0/4),
 \qquad \Psi\ge0\quad\text{on }B_{\widetilde g}(0,r_0).
\end{equation}
The latter inequality uses the admissible endpoint $0$.
Throughout the fixed Euclidean domain,
$\Psi\le C+2\log A$.

We next specify the weighted extension, including its support.
A fixed small initial holomorphic coordinate neighborhood $U$
of $p=\Phi(0)$ lifts to a single sheet through $0$.  Indeed its
short radial paths and their short homotopies lift by local
inversion of $\Phi$, and their $\widetilde g$ lengths equal
their initial target lengths.  The lower metric bound prevents
escape from the coordinate domain.  Uniqueness of local lifts
and the short homotopies make this an inverse branch on $U$.
Choose $U$ uniformly over $p\in K$, with the whole sheet inside
$B_{\widetilde g}(0,r_0/4)$.

In these initial target coordinates $w$, centered at $p$, choose
a linear holomorphic function $\ell$ whose differential has
target norm one and whose pullback differential at $0$ has
Euclidean norm $B$.  Let $\chi$ be a fixed cutoff on this sheet,
equal to one near $0$ and zero near its boundary.  The function
$a=\chi\ell$, extended by zero outside the sheet, is smooth and
uniformly bounded.  Its source $\beta=\dbar a$ is supported in
a fixed target coordinate annulus and satisfies
\[
 |\beta|_{\widetilde g^{-1}}\le C,
 \qquad
 \int_{\supp\beta}dV_{\rm eucl}
 \le c^{-n}\int_{\supp\beta}dV_{\widetilde g}
 \le C.
\]
The last integral is the volume of one initial target annulus,
since the chosen sheet is injective.  These constants do not
contain $B$.

On the same sheet choose a nonpositive logarithmic-pole weight
$\lambda$, equal to $(n+1)\log(|w|^2/\delta^2)$ near $0$ and
cut off to zero within the strict region.  Its negative Hessian
is bounded by $C\widetilde g$ on the cutoff annulus, uniformly
also for smooth pole regularizations.  Choose a fixed large
$k$ so that the weight
\[
 \varphi_*=k\Psi+\lambda+|z|^2
\]
is plurisubharmonic and its Levi form dominates a positive
multiple of $\widetilde g$ on $\supp\beta$.
On that support, $\Psi\ge0$ and $\lambda$ is bounded below.
The scalar weighted $\dbar$ estimate on the fixed Euclidean
ball \cite{Hormander}\cite[Theorem~5.1 and Remark~4.5]{Demailly} gives
$\dbar v=\beta$ with
\begin{equation}\label{chart:weighted-extension}
 \int |v|^2e^{-\varphi_*}\,dV_{\rm eucl}
 \le\int |\beta|^2_{(\dd\varphi_*)^{-1}}
                  e^{-\varphi_*}\,dV_{\rm eucl}\le C.
\end{equation}
For a smooth-weight justification, first convolve $\Psi$ on
a slightly larger Euclidean region and regularize the pole.
For each particular smooth $\widetilde g$, the convolution
scale can be chosen small enough to retain half of
\eqref{chart:strict-envelope} on the compact support regions.
The estimates just obtained are independent of that scale.
Weak $L^2$ limits and Fatou's lemma give
\eqref{chart:weighted-extension} for the singular weight.

The correction $v$ is holomorphic near $0$.  Integrability with
the pole $|w|^{-2(n+1)}$ forces its constant and linear terms
to vanish.  Thus $a-v$ is holomorphic on the Euclidean ball
and has the prescribed differential of norm $B$ at $0$.
Since $\lambda\le0$ and
$\sup\Psi\le C+2\log A$, dropping the weight in
\eqref{chart:weighted-extension} gives
\[
 \|a-v\|_{L^2(dV_{\rm eucl})}^2\le C(1+A)^{2k}.
\]
Interior differentiation of a holomorphic function bounds its
center derivative by this norm, proving
\eqref{chart:static-distortion}.
\end{proof}

\subsection{The volume clock and selected Ricci windows}

The letter $s$ now denotes the volume clock, rather than the time
variable used in the disc envelope.  Set
\begin{equation}\label{chart:clocks}
 \tau=\log t,\qquad s=\rho(o,t),\qquad
 q(s)=\frac{ds}{d\tau}=tR(o,t).
\end{equation}
For all sufficiently large times this is a smooth increasing change
of clock.

We need two kinds of control.  Estimates valid at every sufficiently
late clock will govern the chart transitions and their polynomial
models.  We will also select late times at which the rescaled Ricci
tensor is uniformly pinched at $o$ and the scalar Harnack inequality
approaches equality.  The surrounding intervals provide compactness
for the limits used to contract loops.

\begin{proposition}\label{prop:ricci-windows}
The function $q$ is nondecreasing, $q\ge c>0$ for large $s$, and
\begin{equation}\label{chart:q-growth}
 s\longrightarrow\infty,\qquad
 q\le Ce^\tau,\qquad
 \liminf_{s\to\infty}\frac{q(s)}s\le2.
\end{equation}
Uniformly on every fixed compact set, $\rho(x,t)/s\to1$, and
there is $c_K>0$ such that, eventually on that compact set,
\begin{equation}\label{chart:metric-decay}
 e^{-2s}g\le h_t\le e^{-c_Ks}g.
\end{equation}
There are fixed $a_0\in(0,1)$ and $b_0>0$ such that the least
eigenvalue at $o$ of $t\Ric(h_t)$ relative to $h_t$ satisfies
\begin{equation}\label{chart:ricci-lower}
 \lambda_{\min}(s)\ge b_0q(a_0s)
 \quad\text{for all sufficiently large }s.
\end{equation}
Moreover there are $S_i\to\infty$ and $K_0<\infty$ with
\begin{equation}\label{chart:good-windows}
 q(2S_i)\le K_0q(a_0S_i).
\end{equation}
One can choose $s_i\in[5S_i/4,7S_i/4]$ so that
$(\log q)'(s_i)\to0$.  If $t_i$ corresponds to $s_i$ and
$R_i=R(o,t_i)$, the rescaled flows
\begin{equation}\label{chart:rescaled-flows}
 g_i(z)=R_i h_{t_i+z/R_i}
\end{equation}
have uniform curvature and derivative bounds on each fixed-radius
moving ball about $o$ in a fixed small two-sided time interval.
At $(o,0)$ their scalar curvature is one and their Ricci
eigenvalues have a fixed positive lower bound.
\end{proposition}

\begin{proof}
Proposition~\ref{prop:harnack-completeness} gives monotonicity and
the positive lower bound for $q$, together with $q\le Ce^\tau$.
Hence $s\to\infty$.  If $q(s)>2s$ for every sufficiently large
$s$, integration of $ds/d\tau>2s$ would give
$s\ge ce^{2\tau}$, contradicting the linear-in-$t$ bound for
$\rho(o,t)$.  This proves \eqref{chart:q-growth}.

For a compact family of initial paths from $o$, the power
contraction in Proposition~\ref{prop:harnack-completeness} makes
their $h_t$ lengths tend uniformly to zero.  The factor on the
right of \eqref{chart:integrated-harnack}, with time increment
one, therefore tends uniformly to one on the compact set.  Apply
the inequality in both spatial directions and integrate in time.
More explicitly, for every $\eps>0$ and all sufficiently large
$v$, uniformly on the compact set,
\[
 R(x,v)\le(1+\eps)R(o,v+1),\qquad
 R(o,v)\le(1+\eps)R(x,v+1).
\]
The first inequality also bounds $R(x,v)$ there.  Shifting either
integral endpoint by one costs a bounded amount.  After division
by $s\to\infty$, and then letting $\eps\downarrow0$, the two
inequalities give $\rho(x,t)/s\to1$.  No integrability of the
error factor is required.

Let $\lambda_1,\ldots,\lambda_n$ be the eigenvalues of $h_t$
relative to $g$ at $x$.  They lie in $(0,1]$, and their product
is $e^{-\rho(x,t)}$.  Thus every $\lambda_j\ge e^{-\rho(x,t)}$,
which proves the lower bound of \eqref{chart:metric-decay}.
Choose the instantaneous unitary chart of
Lemma~\ref{chart:local-charts} at $x$.  Its initial-metric
singular lengths are $\lambda_j^{-1/2}$, whose product is
$e^{\rho(x,t)/2}$.  Lemma~\ref{lem:static-distortion} bounds
their largest value by a fixed power of their smallest, uniformly
for $x$ in the compact set.  Since all these lengths are at least
one, their product forces the smallest to be at least
$C^{-1}e^{c\rho(x,t)}$.  This proves the upper bound of
\eqref{chart:metric-decay}, after reducing $c_K$ and increasing
the starting time.

For each fixed nonzero covector $\xi\in T_o^{*(1,0)}M$, the
function
\[
 L_\xi(\tau)=\log|\xi|_{h_{e^\tau}^{-1}}^2
\]
is convex.  Indeed the logarithmic norm of a holomorphic dual
section is plurisubharmonic for a Griffiths-nonnegative tangent
metric, and its restriction to complexified $\tau$ is independent
of the imaginary part.  Its derivative is the Rayleigh quotient
of $t\Ric$ on the dual unit vector.  Normalize $|\xi|_{g^{-1}}=1$.
At $o$, \eqref{chart:metric-decay} gives
\[
 L_\xi(\tau(s))-L_\xi(\tau(a_0s))
 \ge(c_o-2a_0)s.
\]
Choose $a_0<c_o/4$, reducing it further to be less than one if
needed.  Monotonicity of $q$ gives
\[
 \tau(s)-\tau(a_0s)
 =\int_{a_0s}^s\frac{dv}{q(v)}
 \le\frac{s}{q(a_0s)}.
\]
The endpoint derivative of a convex function dominates its secant
slope.  Taking the minimum over all covectors proves
\eqref{chart:ricci-lower}.

The estimates obtained so far hold at every sufficiently late clock.
We now select intervals on which the lower Ricci bound is comparable
to the scalar curvature.  Put $L=2/a_0>1$ and choose $K_0>L^2$.  If
\eqref{chart:good-windows} failed for all large $S$, then
$q(Lr)>K_0q(r)$ for all large $r$.  Iteration and monotonicity
would give $q(r)\ge cr^{\log K_0/\log L}$ for every large $r$,
contradicting \eqref{chart:q-growth}.  Thus such windows occur
arbitrarily far out.  On their middle halves,
\[
 \int_{5S_i/4}^{7S_i/4}(\log q)'(s)\,ds\le\log K_0.
\]
Choose $s_i$ there with
$(\log q)'(s_i)\le2\log K_0/S_i$.

Write $q_i=q(s_i)=t_iR_i$.  On the rescaled clock of
\eqref{chart:rescaled-flows},
\begin{equation}\label{chart:clock-rescaling}
 \frac{ds}{dz}=\frac{q(s)}{q_i+z},\qquad
 \widehat R_i(o,z)=\frac{R(o,t_i+z/R_i)}{R_i}
                 =\frac{q(s)}{q_i+z}.
\end{equation}
Since $q_i\ge c>0$, choose $\delta<c/4$.  As long as $s$ lies
in the good window, $|ds/dz|\le2K_0$ for $|z|\le\delta$.
The chosen $s_i$ lie a distance at least $S_i/4$ from its ends,
so this estimate keeps the whole fixed time interval in the
window for large $i$.  Formula~\eqref{chart:clock-rescaling}
then bounds the rescaled basepoint scalar curvature above and
below by fixed positive constants.  Apply the rescaled
\eqref{chart:integrated-harnack}, using a fixed future-time
buffer, to get a uniform scalar bound on every fixed-radius
moving ball in a smaller interval.  Bisectional nonnegativity
controls $|\Rm|$, and Lemma~\ref{chart:local-charts} supplies
the local higher bounds.  Finally
\eqref{chart:ricci-lower} and \eqref{chart:good-windows} give
at the central point
\[
 \lambda_{\min}(\Ric(g_i))
 \ge\frac{b_0q(a_0s_i)}{q_i}\ge\frac{b_0}{K_0},
 \qquad \widehat R_i(o,0)=1.
\]
\end{proof}

\subsection{Harnack equality and contraction of loops}

The uniform charts are still only local biholomorphisms.  To continue
their inverse germs to single-valued maps on compact subsets of $M$,
we must remove possible monodromy.  The selected windows do this by
producing equality in a limiting Harnack inequality.  We will show
that this equality forces real strict concavity of the scalar
curvature, then use its gradient flow to contract loops in the
corresponding small intrinsic balls.

\begin{proposition}\label{prop:simple-connectivity}
The manifold $M$ is simply connected.
\end{proposition}

\begin{proof}
Take the rescaled flows in Proposition~\ref{prop:ricci-windows}.
Their pointed pullbacks on tangent balls have a smooth local
subsequential limit, denoted $g(z)$.  Pass to a further subsequence
so that
\[
 k_i(z)=\frac1{q_i+z}\longrightarrow k(z),\qquad k_z=-k^2.
\]
This includes $k=0$ when $q_i\to\infty$.  Denote the limit
scalar curvature by $R$ and its Ricci form by $A$.  At the
central spacetime point, $R=1$ and $A\ge b\,g$ with $b>0$.
Formula~\eqref{chart:clock-rescaling} gives the exact identity
\[
 \left[\partial_z\widehat R_i+k_i\widehat R_i\right]_{(o,0)}
       =(\log q)'(s_i)\longrightarrow0.
\]
The limit therefore has $S:=R_z+kR=0$ at the central point.

The matrix Harnack inequality has the rescaled form
$A_z+A g^{-1}A+kA\ge0$.  Its trace is $S$, because
$R_z=\tr_g A_z+|A|_g^2$.  Thus this nonnegative matrix vanishes
at the point in question.  Also $A_z=\dd R$ under K\"ahler--Ricci
flow.  If $H_{j\bar l}=\nabla_j\nabla_{\bar l}R$, then
\begin{equation}\label{chart:equality-mixed-hessian}
 H=-A^2-kA
\end{equation}
there, writing endomorphism products in a unitary frame.
On a small spacetime neighborhood, $A$ is positive definite.
Minimizing the scalar Harnack inequality over $X$ gives
\begin{equation}\label{chart:equality-Z}
 Z=S-Q\ge0,\qquad Q=|\partial R|_{A^{-1}}^2.
\end{equation}
At the equality point $S=0$, so $\partial R=0$ and $Z=0$.

We give the scalar Hessian calculation, since a vanishing
Harnack quantity by itself does not supply real strict concavity.
Use $R_z=\Delta R+|A|^2$, $A_z=\dd R$, and $k_z=-k^2$.
The derivative of the Laplacian on a scalar is
$(\partial_z\Delta)R=\langle A,\dd R\rangle$, while
\[
 \partial_z|A|^2=2\langle A,\dd R\rangle+2\tr(A^3).
\]
Consequently
\begin{equation}\label{chart:equality-evolution}
 (\partial_z-\Delta)S
 =3\langle A,\dd R\rangle+2\tr(A^3)
     +k|A|^2-k^2R.
\end{equation}
At the equality point, substitution of
\eqref{chart:equality-mixed-hessian} reduces this to
$-\tr\bigl(A(A+k\Id)^2\bigr)$.
Diagonalize $A$ there, with eigenvalues $\lambda_j>0$.
Since $\partial R=0$, derivatives of the coefficients of
$A^{-1}$ make no contribution to $\Delta Q$ or $Q_z$, and
\[
 Q_z=0,\qquad
 \Delta Q=\sum_{j,l}\lambda_j^{-1}
       \left(|R_{j\bar l}|^2+|R_{jl}|^2\right).
\]
The first square sum is exactly
$\tr\bigl(A(A+k\Id)^2\bigr)$ by
\eqref{chart:equality-mixed-hessian}.  Hence at this point
\[
 (\partial_z-\Delta)Z
 =\sum_{j,l}\lambda_j^{-1}|R_{jl}|^2\ge0.
\]
On the other hand $Z\ge0$ has an interior spacetime zero, so
$Z_z=0$ and $\Delta Z\ge0$ there.  The displayed sum must
vanish.  Thus the pure complex Hessian $R_{jl}$ is zero, while
the mixed Hessian is negative definite by
\eqref{chart:equality-mixed-hessian}.  The real Hessian of $R$
is therefore negative definite at the center, with a definite
bound furnished by the positive lower bound for $A$.

This conclusion holds uniformly on actual small metric balls
in the original rescaled manifolds, not only on their immersed
parametrizations.  Smooth convergence gives the center Hessian
bound and $|\nabla\widehat R_i|(o,0)\to0$.  The intrinsic
third-derivative estimates from
Proposition~\ref{prop:ricci-windows} allow one to parallel
transport the Hessian along any minimizing radial geodesic.
Shrinking a fixed radius $r_*>0$, independently of $i$, gives
constants $a>0$ with
\begin{equation}\label{chart:intrinsic-concavity}
 \operatorname{Hess}_{g_i(0)}\widehat R_i
       \le-a g_i(0)
 \quad\text{on }B_{g_i(0)}(o,2r_*),\qquad
 |\nabla\widehat R_i|(o,0)\longrightarrow0.
\end{equation}
Indeed the change in the transported Hessian is bounded by
$Cr_*$, so one chooses $r_*$ after the uniform center bound.
This argument uses intrinsic derivative estimates along actual
geodesics and is unaffected by multiple sheets of a local chart.

Hold the complete metric $g_i(0)$ fixed and flow by the real
vector field $Y_i=\nabla\widehat R_i(\cdot,0)$.  Along a
minimizing radial geodesic of length $r\le2r_*$,
\[
 \langle Y_i,\partial_r\rangle
 \le |Y_i(o)|-ar.
\]
For large $i$ this is negative on $r=r_*$.  At cut points,
first variation along any minimizing radial segment gives the
same upper bound for the upper directional derivative of
distance.  Thus the closed ball of radius $r_*$ traps trajectories
starting inside it.  Completeness makes this closed ball compact,
so these trajectories exist for all positive flow time.
For a carried tangent vector $V$, \eqref{chart:intrinsic-concavity}
gives
\[
 \frac{d}{dv}|V|_{g_i(0)}^2
 =2\operatorname{Hess}\widehat R_i(V,V)
 \le-2a|V|_{g_i(0)}^2.
\]
Every loop contained well inside the trapping ball consequently
has carried length tending exponentially to zero.

Let a fixed piecewise smooth loop in $M$ be given.  Its image
and a finite collection of paths from $o$ lie in a fixed compact
set.  Their $h_{t_i}$ lengths tend to zero by
\eqref{chart:metric-decay}, and $R_i\le C$ by
Proposition~\ref{prop:harnack-completeness}.  Thus this loop
lies inside $B_{g_i(0)}(o,r_*/4)$ for sufficiently large $i$.
Fix one such $i$.  The compact trapping ball has a positive
minimum injectivity radius for this one smooth metric.  A
sufficiently short carried loop lies in a normal ball and is
contractible.  The gradient flow up to that finite time is a
homotopy from the original loop to it.  Approximation of
continuous loops by piecewise smooth loops finishes the proof.
\end{proof}

\subsection{Inverse charts and one contraction exponent}

Simple connectivity permits continuation of the inverse germs, but
their first contraction estimates will depend on the compact set.
The following Liouville principle, applied to limits of logarithmic
norms, will make the exponent common to all compact sets.  We state it
in the required regularity class; it is also a special case of
\cite[Theorem~9]{LiuThreeCircle}.

\begin{lemma}\label{chart:psh-liouville}
Every bounded-above plurisubharmonic function on $M$ is constant,
including a function which is not assumed continuous.
\end{lemma}

\begin{proof}
Allow the constant $-\infty$, and otherwise let $v$ be such a
function.  The compact-ball maximum
$m(r)=\max_{\overline B_g(o,r)}v$ is finite and nondecreasing.
It is convex as a function of $\log r$.  To see this, fix
$0<r_1<r_2$ and compare on the closed annulus with the
logarithmic chord
\[
 H(r)=m(r_1)
  +\frac{m(r_2)-m(r_1)}{\log r_2-\log r_1}
      (\log r-\log r_1).
\]
If its slope is zero the desired comparison is immediate from
the definition of $m(r_2)$.  For positive slope, a positive
maximum of $v-H(\dist_g(o,\cdot))$ would occur in the annulus
interior.  At such a point choose a minimizing radial segment.
The length index form with linear parallel test fields in the
$J$-radial direction gives an upper Hessian bound for distance.
More explicitly, let $T$ be the final radial unit vector, let
$T(\ell)$ denote its parallel extension, and let $\bar r$ be
the length upper support obtained from these trial paths.  At
the endpoint of a segment of length $r$ the real Hessian satisfies
\[
 \operatorname{Hess}\log\bar r(T,T)
 +\operatorname{Hess}\log\bar r(JT,JT)
 \le-\frac1{r^3}\int_0^r
       \ell^2\Rm(T(\ell),JT(\ell),JT(\ell),T(\ell))\,d\ell<0.
\]
Indeed the radial term is $-1/r^2$, while the linear test field
in the $JT$ direction gives $1/r^2$ minus the displayed
curvature integral.  This is the strictly negative Levi
derivative on the radial complex line.

This support argument applies also at cut points: smoothly vary
the endpoint and use trial paths along the chosen minimizing
segment.  Their lengths give a smooth upper support with the
same index-form bound.  Since $H$ is increasing, it produces
a smooth upper test for $v$ at the hypothetical positive
maximum, with a negative Levi derivative.  A plurisubharmonic
function admits no such upper test, whether or not it is
continuous.  Hence $v\le H$ on the annulus, proving the
claimed convexity.  A nondecreasing convex function of
$\log r\in\R$ which is bounded above must be constant.
Upper semicontinuity implies
$\lim_{r\downarrow0}m(r)=v(o)$, so $v$ attains its global
maximum at $o$.  The plurisubharmonic maximum principle on
connected $M$ shows that $v$ is constant.
\end{proof}

Write $J_dF$ for the Taylor jet of $F$ through total degree $d$ at zero,
and fix a coefficient norm on each finite-dimensional jet space.

\begin{proposition}\label{prop:shrinking-charts}
For every sufficiently large volume clock $s$, there is a
holomorphic local biholomorphism
$\Phi_s:B_{\C^n}(r_0)\to M$ centered at $o$, with unitary
center derivative for the metric at clock $s$, and
\begin{equation}\label{chart:uniform-chart-metric}
 C^{-1}g_{\rm eucl}\le\Phi_s^*h_{t(s)}\le Cg_{\rm eucl}.
\end{equation}
For fixed $0<r_1<r_0$, chosen sufficiently small, the transition
germs $F_{s,u}=\Phi_u^{-1}\Phi_s$, $u\ge s$, extend to
nonsingular holomorphic maps on $B(r_1)$ with a common bound.
They fix $0$, compose as germs, and satisfy
\begin{equation}\label{chart:center-transitions}
 \|F_{s,u}'(0)\|\le1,\qquad
 |\det F_{s,u}'(0)|=e^{-(u-s)/2},\qquad
 \|F_{s,u}'(0)^{-1}\|\le e^{(u-s)/2}.
\end{equation}
There are coherent inverse branches $f_s$ on neighborhoods of
any fixed compact exhaustion set for all sufficiently large
$s$, satisfying $f_s(o)=0$ and $\Phi_sf_s=\Id$.  A constant
$\sigma>0$, common to all compact sets, satisfies
\begin{equation}\label{chart:common-exponent}
 \sup_K|f_s|\le e^{-\sigma s}
 \quad\text{for every fixed compact }K
 \text{ and all sufficiently large }s.
\end{equation}
For every fixed jet order $d$ there is $C_d$ such that
\begin{equation}\label{chart:transition-jets}
 |J_dF_{s,u}|\le C_d\min\{1,e^{C_ds-\sigma u}\}
 \qquad(u\ge s\text{ sufficiently large}).
\end{equation}
There are $c,\gamma_0>0$ such that, for each fixed
$0<\gamma\le\gamma_0$, one has
\begin{equation}\label{chart:transition-small-balls}
 \sup_{|z|\le e^{-\gamma s}}|F_{s,u}(z)|
 \le e^{-c\gamma u}
 \qquad\text{for every }u\ge s\ge s_\gamma.
\end{equation}
The constants in the last two estimates are independent of the
later endpoint $u$; the starting threshold in
\eqref{chart:common-exponent} may depend on $K$.
\end{proposition}

\begin{proof}
Take the charts of Lemma~\ref{chart:local-charts} at $o$.
We use the elementary short-path lifting property of these
possibly noninjective charts; compare
\cite[Lemmas~2.2--2.3 and Corollary~2.2]{ChauTamStein}.
A path beginning at the center and having sufficiently small
current-metric length lifts from $0$ by local inversion.
The lower bound in \eqref{chart:uniform-chart-metric} bounds
the Euclidean length of its lift.  If the given length is
smaller than the fixed distance to the chart boundary, the
lift cannot escape, and hence extends over the whole path.
The same reasoning applies to a homotopy all of whose paths
obey this length bound.  Local uniqueness of lifts supplies
uniqueness throughout the homotopy.

For $u\ge s$, the image under $\Phi_s$ of a short radial path
has $h_{t(u)}$ length at most its $h_{t(s)}$ length.  Lift it
through $\Phi_u$ from $0$.  Short radial homotopies give a
well-defined map on a fixed ball $B(r_1)$; locally this map
is $\Phi_u^{-1}\Phi_s$, so it is holomorphic with nonsingular
derivative.  Its lift has a fixed Euclidean length bound, and
therefore the maps are uniformly bounded on $B(r_1)$.
They compose wherever both sides are defined near $0$ by
uniqueness of local inverses.  This is the local transition
construction of \cite[Lemmas~3.1--3.2]{ChauTamStein} with the
metric bounds already established here.

Unitary center normalization and $h_{t(u)}\le h_{t(s)}$
give the first estimate in \eqref{chart:center-transitions}.
Taking determinants of these normalizations and using
$\det h_{t(s)}(o)=e^{-s}\det g(o)$ gives the determinant
identity.  All singular values are at most one and their
product is $e^{-(u-s)/2}$, so the least is at least this
product.  This proves the inverse estimate.

We now use simple connectivity to continue the inverse germs.
Fix a connected relatively compact neighborhood of a compact
set and $o$.  Cover its closure by finitely many small convex
initial target neighborhoods with connected pairwise
intersections, enlarging the connected neighborhood if
necessary to include fixed paths from $o$ to their centers.
Choose one such path to each center.  On each target
neighborhood it extends the inverse germ at $o$ by path
lifting.  For each nonempty overlap, use one point in the
overlap to form the two-path loop.  By
Proposition~\ref{prop:simple-connectivity} this loop bounds
a fixed piecewise smooth homotopy.  There are only finitely
many paths and homotopies; their images lie in a fixed
compact set and their initial lengths have a uniform bound.
Equation~\eqref{chart:metric-decay} makes every path in these
homotopies sufficiently short in the current metric when
$s$ is large.  All lifts stay inside the chart.  Thus the
two inverse germs agree at the chosen overlap point, and
then throughout the connected overlap by analytic
continuation.  They patch to an inverse $f_s$ on a
neighborhood of the compact set.

Choose nested connected relatively compact open exhaustion
sets $D_m$, and choose nondecreasing thresholds $T_m\ge m$
so that all shortness conditions for $D_m$ hold whenever
$s\ge T_m$.  At clock $s$, use the construction on
$D_{\max\{m:T_m\le s\}}$.  On a smaller set, the branches
obtained from two larger constructions agree by continuation
from $o$.  This gives the stated coherent branches with
domains exhausting $M$.  The length estimate for a lift,
together with \eqref{chart:metric-decay} on the finite path
compact, gives for each fixed $K$
\begin{equation}\label{chart:local-exponent}
 \sup_K|f_s|\le C_Ke^{-a_Ks}
\end{equation}
with $a_K>0$.  It also gives
$f_u=F_{s,u}\circ f_s$ on any fixed compact set whenever
$s$ is sufficiently large and $u\ge s$: both sides are
the same continued inverse, and
\eqref{chart:local-exponent} keeps $f_s$ in the transition
domain.

It remains to make the exponent independent of the compact.
The functions $v_s=s^{-1}\log|f_s|$ are plurisubharmonic on
their exhausting domains and locally bounded above.
For any sequence $s\to\infty$,
Lemma~\ref{lem:psh-compactness-bridge} and diagonal extraction
give either convergence to
$-\infty$ locally uniformly from above, or local $L^1$
convergence to a plurisubharmonic function on $M$.
In the second case the upper-regularized limit is at most
zero everywhere by \eqref{chart:local-exponent}.
Lemma~\ref{chart:psh-liouville} makes it constant.  On one
fixed neighborhood of $o$, \eqref{chart:local-exponent}
bounds this constant by a fixed negative number, say
$-a_*$.  The constant-limit conclusion of
Lemma~\ref{lem:psh-compactness-bridge} then gives
$\limsup\sup_Kv_s\le-a_*$ on every fixed compact along
this subsequence.  The collapsing case gives the same
conclusion with $-\infty$.  A sequence violating
\eqref{chart:common-exponent}, with $\sigma=a_*/2$, would
contradict these alternatives.  This proves a common
exponent on the full family of clocks.

For the two remaining transition estimates, transfer the common
exponent from a fixed initial neighborhood back to a small ball in
the chart domain.  Choose a fixed compact initial neighborhood $K_0$
of $o$ with positive initial distance from $o$ to its
boundary.  Until a path $\Phi_s(rz)$ first exits $K_0$,
\eqref{chart:metric-decay} and
\eqref{chart:uniform-chart-metric} bound its initial length
by $Ce^s|z|$.  It follows that, for a fixed $C_0>1$ and
all large $s$,
\[
 \Phi_s\bigl(B(e^{-C_0s})\bigr)\subset K_0.
\]
For every $u\ge s$, the map $f_u\Phi_s$ on this ball is
the same local inverse composition as $F_{s,u}$ near zero,
and hence everywhere there.  Equation~\eqref{chart:common-exponent}
therefore gives the bound
\begin{equation}\label{chart:transition-inner-bound}
 \sup_{B(e^{-C_0s})}|F_{s,u}|\le e^{-\sigma u}
 \qquad\text{for all }u\ge s\text{ and large }s.
\end{equation}
Cauchy's estimates on this inner ball give the factor
$e^{dC_0s-\sigma u}$ for derivatives of order at most $d$.
Cauchy's estimates on a fixed transition ball give a uniform
bound.  Combining the two proves
\eqref{chart:transition-jets}.

For completeness, fix an outer radius $R<r_1$ and a common
bound $B_*$ for the transitions there.  Euclidean
three-circle interpolation, applied to their scalar linear
projections and then taking the supremum, gives at radius
$e^{-\gamma s}$, for fixed small $\gamma>0$,
\[
 \log\sup_{B(e^{-\gamma s})}|F_{s,u}|
 \le-\theta_s\sigma u+(1-\theta_s)\log B_*,
 \qquad
 \theta_s=\frac{\log(R/e^{-\gamma s})}
                 {\log(R/e^{-C_0s})}
 \ge\frac{\gamma}{2C_0}
\]
for all sufficiently large $s$.  In this formula choose
$\gamma_0<C_0$, so the middle radius lies between the
inner and outer radii.  Since $u\ge s$, the fixed term
$\log B_*$ is absorbed uniformly for every later endpoint,
giving \eqref{chart:transition-small-balls} with, for
example, $c=\sigma/(4C_0)$ and a threshold depending on
$\gamma$.  This proves all the stated uniformities.
\end{proof}
